% Generated draft: requires mathematical and source review.
\documentclass[11pt]{article}
\usepackage[margin=25mm]{geometry}
\usepackage{fontspec}
\setmainfont{Latin Modern Roman}
\usepackage{amsmath,amssymb,mathtools,unicode-math}
\setmathfont{Latin Modern Math}
\usepackage{enumitem,needspace,xurl,hyperref}
\setlength{\parindent}{0pt}
\setlength{\parskip}{5pt}
\setlength{\emergencystretch}{3em}
\newcommand{\N}{\mathbb{N}}
\newcommand{\Z}{\mathbb{Z}}
\newcommand{\Q}{\mathbb{Q}}
\newcommand{\R}{\mathbb{R}}
\newcommand{\C}{\mathbb{C}}
\newcommand{\F}{\mathbb{F}}
\newcommand{\A}{\mathbb{A}}
\newcommand{\PP}{\mathbb{P}}
\newcommand{\E}{\mathbb{E}}
\newcommand{\eps}{\varepsilon}
\newcommand{\dd}{\,\mathrm{d}}
\begin{document}
\section*{UnsolvedMath ID: 20003358}
\textbf{Source problem code:} catalogue entry
\subsection{Definitions and mathematical statement}
The extracted ``Conjecture 29'' is a section-boundary continuation, rather than a new conjecture. The preceding arithmeticity setting, recovered from neighboring local records 20003355--20003357, is
\[
G=\mathrm{SL}_2(\R)^k,\quad k\geq2,\qquad
u_+(x)=\left(\begin{pmatrix}1&x_i\\0&1\end{pmatrix}\right)_{i=1}^k,
\quad
u_-(y)=\left(\begin{pmatrix}1&0\\y_i&1\end{pmatrix}\right)_{i=1}^k,
\]
with \(U^\pm=\nu_\pm(\R^k)\). A full lattice \(L\subset\R^k\) is irreducible here if each coordinate projection is injective on \(L\), equivalently, every nonzero lattice vector has all coordinates nonzero. This is the source's specific meaning of irreducibility.

Let \(\mathcal L_k=\mathrm{GL}_k(\R)/\mathrm{GL}_k(\Z)\) be the space of full, unoriented lattices, equipped locally with a smooth Riemannian metric, and set
\[
\mathcal E_k=\{L\in\mathcal L_k:L\text{ irreducible and there exists an irreducible full lattice }M\subset\R^k
\text{ for which }\langle\nu_+(L),\nu_-(M)\rangle\text{ is discrete in }G\}.
\]
The substantive dimension assertion reported by the supplied research aid, for \(k\geq3\), is
\[
\dim_{\mathrm H}\mathcal E_k=k.
\]
The AIM text attributes this assertion to reference [152]; the aid identifies the result with Oh. The assertion concerns the full lattice space, including covolume. Its covolume-one version, as reported in the aid, is
\[
\mathcal L_k^{(1)}=\{L:\operatorname{covol}(L)=1\},\qquad
\dim_{\mathrm H}(\mathcal E_k\cap\mathcal L_k^{(1)})=k-1.
\]
Positive scalar dilation preserves eligibility, so these are different normalizations of the same geometric locus, not identical dimension statements. The initial setup allows \(k=2\), but the supplied aid states the dimension result for \(k\geq3\); this draft does not infer its asserted range at \(k=2\).

For clarity about the aid's finite-index qualification, let \(\mathcal F_k\) consist of lattices \(\nu_+^{-1}(\Delta\cap U^+)\), where \(\Delta\leq G\) is discrete and both \(\Delta\cap U^\pm\) are irreducible lattices in \(U^\pm\). The reported passage from full intersections to the exact input-lattice locus is
\[
\mathcal E_k=\{L:\exists L'\in\mathcal F_k,\ L\subset L',\ [L':L]<\infty\}.
\]
Thus a theorem formulated for full horospherical intersections requires a finite-index correspondence argument before it applies to arbitrary input lattices. The aid describes these as countably many finite Hecke correspondences and reports preservation of the relevant transverse upper box dimension zero property, a positive-diagonal orbit argument, and covolume normalization. It also reports that irreducibility is dense and conull in \(\mathcal L_k\). These are supplied research claims requiring independent verification, not mathematical review performed here; precise charts, measures, and the cited rigidity theorem must be checked in their original sources.

The remainder starts Section 11, ``Symbolic coding.'' It is background: a geometric code records successive fundamental-domain sides met by an oriented geodesic, while an arithmetic code uses continued-fraction expansions of its endpoint coordinates in an appropriate hyperbolic model. Boundary hits, nonunique expansions, and the chosen domain or cross-section require conventions. No new coding conjecture or equivalence of these two codes is stated in this fragment.
\subsection{Short English statement}
The fragment reports a dimension theorem for irreducible upper unipotent lattices that can be paired with an opposite lattice to generate a discrete group. The reported dimension is \(k\) before fixing covolume and \(k-1\) after normalization, for \(k\geq3\). The appended discussion of geodesic coding belongs to the next section.
\subsection{Sources}
\begin{itemize}
\item Local UnsolvedMath registry, permanent ID 20003358, catalogue entry; canonical location \texttt{aim-other-notes.json}, notes[164]. Neighboring local records 20003355--20003357 supply the preceding definitions of \(G\), \(U^\pm\), and irreducibility.
\item American Institute of Mathematics, \emph{Emerging applications of measure rigidity}, end of Section 10 and beginning of Section 11: \url{https://aimath.org/WWN/measrigid/measrigid.pdf}. The dimension sentence cites [152], identified in the supplied aid with Oh; its complete bibliographic entry is not supplied.
\item The preceding arithmeticity conjecture is attributed to G. Margulis and A. Selberg and was communicated by H. Oh; the intervening implication cites [151]. These are context, not a new attribution of the dimension theorem.
\item The supplied aid invokes an Einsiedler--Katok--Lindenstrauss full-intersection result and finite Hecke correspondences, but provides no full bibliographic entry or theorem number for the dimension transfer.
\item The symbolic-coding background cites [27], [30, 173, 174, 175], [1, 2], and [100, 75, 101, 102], and credits early constructions to M. Morse, E. Artin, and G. Hedlund. Full bibliography entries are absent from this record.
\end{itemize}
\end{document}
